\noindent
\begin{minipage}[t]{0.48\textwidth}
    \begin{center}
        BỘ GIÁO DỤC VÀ ĐÀO TẠO\\
        \textbf{ĐẠI HỌC PHENIKAA}
    \end{center}
\end{minipage}
\hfill
\begin{minipage}[t]{0.48\textwidth}
    \begin{center}
        {\fontsize{8pt}{11pt}\selectfont \textbf{CỘNG HÒA XÃ HỘI CHỦ NGHĨA VIỆT NAM}}\\
        \textbf{Độc lập - Tự do - Hạnh phúc}
    \end{center}
\end{minipage}

\vspace{0.5cm}

% ====== TIÊU ĐỀ CHÍNH ======
\begin{center}
    \textbf{\Large NHẬN XÉT ĐỒ ÁN TỐT NGHIỆP CỦA GIẢNG VIÊN PHẢN BIỆN}
\end{center}

\vspace{0.5cm}

% ====== THÔNG TIN SINH VIÊN ======
\noindent
\textbf{Giảng viên phản biện:} Phạm Trung Dũng\\
\textbf{Khoa:} Khoa học máy tính\\
\textbf{Tên đề tài:} \textit{Xây dựng mô hình cho bài toán người giao hàng kết hợp phương tiện bay không người lái}\\
\textbf{Sinh viên thực hiện:} Lê Hữu Trung \hfill \textbf{Lớp:} K13--CNTT--KHDL\\
\textbf{Giảng viên hướng dẫn:} TS. Dương Thị Kim Huyền

\vspace{0.3cm}

% ====== NỘI DUNG NHẬN XÉT ======
\begin{center}
    \textbf{NỘI DUNG NHẬN XÉT}
\end{center}

\vspace{0.3cm}
\noindent\textbf{I. Nhận xét ĐA/KLTN:}

\begin{itemize}
    \item[-] Bố cục, hình thức trình bày:\\
    Báo cáo được trình bày đúng và đầy đủ theo hình thức của một báo cáo đồ án hoàn chỉnh. Báo cáo có cấu trúc bố cục gồm 3 chương. Chương 1 trình bày giới thiệu về bài toán và các nghiên cứu liên quan. Chương 2 trình bày mô hình MILP cho bài toán giao hàng. Chương 3 trình bày thực nghiệm và kết quả. Hình ảnh minh họa và biểu đồ được chèn hợp lý, hỗ trợ tốt cho việc diễn giải nội dung.

    \item[-] Đảm bảo tính cấp thiết, hiện đại, không trùng lặp của đề tài:\\
    Đề tài giải quyết bài toán logistics, tối ưu việc giao hàng. Nội dung có tính ứng dụng cao, giải quyết được các bất cập trong phương pháp giao hàng truyền thống. Đề tài không trùng lặp với các đồ án đã biết trước đó.

    \item[-] Nội dung của đề tài:\\
    Nội dung đồ án bao quát đầy đủ các khía cạnh từ phát biểu bài toán, các nghiên cứu liên quan, phương pháp giải quyết bài toán, thực nghiệm và kết luận. Phần lý thuyết được trình bày rõ ràng, trình bày thuật toán khoa học.

    \item[-] Mức độ thực hiện:\\
    Sinh viên kết hợp giữa bài toán tối ưu lộ trình truyền thống và điều phối drone. Mô hình được xây dựng có cơ sở lý thuyết rõ ràng, phản ánh các ràng buộc thực tế như quãng đường, tải trọng và thời gian. Sinh viên đã triển khai thử nghiệm mô phỏng, cho kết quả khả quan về giảm thời gian và chi phí giao hàng.
\end{itemize}

% \vspace{0.3cm}
\noindent\textbf{II. Kết quả đạt được:}\\[0.2cm]
Sinh viên đã xây dựng thành công mô hình cho bài toán người giao hàng kết hợp phương tiện bay không người lái, thể hiện sự hiểu biết về tối ưu hóa và điều phối đa phương tiện. Mô hình đề xuất phản ánh được các ràng buộc thực tế và được kiểm chứng qua các bộ dữ liệu mô phỏng. Kết quả cho thấy mô hình giúp giảm đáng kể thời gian và chi phí giao hàng so với phương pháp truyền thống. Sản phẩm đạt được có ý nghĩa thực tiễn, mở hướng cho các nghiên cứu và ứng dụng trong lĩnh vực logistics thông minh.

\vspace{0.3cm}
\noindent\textbf{III. Ưu nhược điểm:}

\begin{itemize}

\item Ưu điểm:
    \begin{itemize}
        \item Đề tài thực tiễn, một hướng đi khả thi đang được phát triển là dùng drone để giao hàng.
        \item Mô hình toán học chặt chẽ – Đồ án xây dựng hai mô hình MILP (5-index và 3-index) giúp mô tả chi tiết và tối ưu hóa bài toán, thể hiện năng lực tư duy toán học tốt.
        \item Thực nghiệm đa dạng, có cơ sở – Tiến hành thử nghiệm trên nhiều bộ dữ liệu chuẩn với công cụ CPLEX, kết quả được trình bày rõ ràng, có số liệu minh chứng cụ thể.
    \end{itemize}

\item Nhược điểm:
    \begin{itemize}
        \item Giới hạn quy mô bài toán, mô hình hiện tại chỉ giải được các bộ dữ liệu nhỏ ($\leq 20$ đỉnh); khi mở rộng quy mô, thời gian và bộ nhớ tăng theo cấp số mũ.
        \item Chưa xét đến các thông số thực tiễn như chi phí vận hành, giới hạn tải drone, điều kiện môi trường hoặc khung thời gian giao hàng, khiến kết quả chưa hoàn toàn phản ánh thực tế.
    \end{itemize}
\end{itemize}


% \vspace{0.3cm}
\noindent\textbf{IV. Kết luận:}\\[0.2cm]
\ding{51} \textbf{Đồng ý cho bảo vệ} \hspace{2.5cm} \ding{113} \textbf{Không đồng ý cho bảo vệ}

\vspace{0.3cm}
\noindent
\begin{minipage}[t]{0.48\textwidth}
\end{minipage}
\hfill
\begin{minipage}[t]{0.48\textwidth}
    \begin{center}
        Hà Nội, ngày 13 tháng 10 năm 2025\\[0.2cm]
        \textbf{GIẢNG VIÊN PHẢN BIỆN}\\[0.2cm]
        (Ký, ghi rõ họ tên)\\[1.2cm]
        \textbf{Phạm Trung Dũng}
    \end{center}
\end{minipage}